\documentclass[10pt,a4paper]{amsart}
\usepackage[margin=19mm]{geometry}
\usepackage{amsmath,amssymb,mathtools}
\usepackage[scr=rsfs,cal=euler]{mathalfa}
\usepackage{xfrac}
\usepackage[unicode,hidelinks,bookmarks=false]{hyperref}
\newcommand{\ie}{i.e.\ }
\newcommand{\cf}{cf.\ }
\newcommand{\resp}{respectively\ }
\newcommand{\Subsection}{\subsection}
\providecommand{\nobreakdash}{\nobreak\hbox{-}}
\setlength{\emergencystretch}{2em}
\multlinegap0pt
\usepackage{bm}
\usepackage{bbm}
\usepackage{dsfont}
\usepackage{xspace}
\usepackage{cancel}
\usepackage{mathtools}
\usepackage{amsthm}
\makeatletter
\@ifundefined{lemma}{\newtheorem{lemma}{Lemma}}{}
\@ifundefined{proposition}{\newtheorem{proposition}{Proposition}}{}
\makeatother
\newcommand{\nn}{\ensuremath{\mathbb{N}}}
\title[Bruhat Intervals in the Infinite Symmetric Group are Cohen-M]{Bruhat Intervals in the Infinite Symmetric Group are Cohen-Macaulay}
\author{Nathaniel Gallup, Leo Gray}
\date{Source: arXiv:2601.12195v1 (CC BY 4.0)}
\begin{document}
\maketitle

\noindent\emph{Source and licence.} Bruhat Intervals in the Infinite Symmetric Group are Cohen-Macaulay. Version arXiv:2601.12195v1 (CC BY 4.0), distributed under the Creative Commons Attribution 4.0 International licence, as recorded in the arXiv licence metadata for this exact version and shown on the abstract page https://arxiv.org/abs/2601.12195v1. Full rights notice: licenses/arxiv-2601.12195-CC-BY-4.0.txt.\par\smallskip

\noindent\textit{Editorial scope.} Counted unit: the Proposition whose source label is \texttt{prop: cover relations in the bruhat order} in the article (source0.tex lines 546--552, source label \texttt{prop: cover relations in the bruhat order}), delivered together with the article's own complete original proof of it (source0.tex lines 554--588, one \texttt{proof} environment of 35 lines). No mathematics is written, repaired or re-derived: statement and proof are the article's own text line for line. Required context retained uncounted: prop: region criterion for bruhat order (lines 446--454); lem: ones in boxes after applying a transposition (lines 473--478); prop: transposition going up in bruhat (lines 506--511); lem: criteria for equal permutation matrices (lines 529--537). Every definition, display and label of those lines is retained unchanged. Omitted from this package and not redistributed: 1--445, 455--472, 479--505, 512--528, 538--545, 553--553, 589--914. Nothing inside a retained excerpt is omitted or altered: every retained body is delivered exactly as the article's lines are written, so no omission receipt of any kind accompanies this package. Citations: the retained excerpts carry no citation command at all, so no bibliography entry of the article is transcribed and no reference is redistributed. Reference adaptation: none was needed and none was made; every reference inside the delivered unit resolves inside it, so no unresolved reference or citation is produced. Printed slips kept literal: no printed word, symbol, hypothesis or number of the counted statement or of its proof was altered. Layout and class: the article's own class is \texttt{amsart} with packages including blkarray, nicematrix, amssymb, graphicx, caption, subcaption, amsmath, bbm, amsthm, amsfonts, setspace, mathrsfs, hyperref, tensor; the wrapper is the lane's pinned \texttt{amsart} editorial wrapper at 10pt on A4, which restates the theorem-like environments the retained text needs and adds the article's own macro definitions that the retained text needs (\texttt{\textbackslash nn}), with extra wrapper packages (bm, bbm, dsfont, xspace, cancel, mathtools). The delivered numbering is the unit's own. Assets: no retained span contains a figure environment, a graphics-inclusion command, a PDF-page inclusion, a table or a TikZ picture; no image file is copied into this package, the package declares no asset dependency and no image file is redistributed with it. Licence: version arXiv:2601.12195v1 of this article is distributed under the Creative Commons Attribution 4.0 International licence, as recorded in the arXiv licence metadata for this exact version and shown on the abstract page https://arxiv.org/abs/2601.12195v1. A full rights notice is delivered at licenses/arxiv-2601.12195-CC-BY-4.0.txt. Characters: every retained excerpt is pure ASCII, so the delivered engine, XeLaTeX with its default Latin Modern text font, reports no missing character.
\begin{proposition}\label{prop: region criterion for bruhat order}
Given $\sigma, \omega \in S_\infty$, the following are equivalent. 

\begin{enumerate}
    \item $\sigma \leq_\text{Bruhat} \omega$
    \item $r_{a , b}(\sigma) \geq r_{a , b}(\omega)$ for all $a , b \in \nn$
    \item $r_{\geq a , b}(\sigma) \leq r_{\geq a , b}(\omega)$ for all $a , b \in \nn$
\end{enumerate}
\end{proposition}

\begin{lemma}\label{lem: ones in boxes after applying a transposition}
    Suppose $\sigma \in S_\infty$ and $p , q \in \nn$ are such that $p < q$ and $\sigma(p) < \sigma(q)$. Then we have the following formulas. 
    \begin{equation*}
        r_{a , b}(\sigma) = \left\{ \begin{matrix} r_{a , b}(\sigma \circ (p , q)) & \text{ if } & b < p, b \geq q , a < \sigma(p), \text{ or } a \geq \sigma(q) \\ r_{a , b}(\sigma \circ (p , q)) + 1 & \text{ if } & \sigma(p) \leq a < \sigma(q) \text{ and } p \leq b < q  \end{matrix} \right. 
    \end{equation*} 
\end{lemma}

\begin{proposition}\label{prop: transposition going up in bruhat}
    For any $\sigma \in S_\infty$, and any $p < q$, we have that $\sigma <_\text{Bruhat} \sigma \circ (p , q)$ if and only if the following condition is satisfied.
    \begin{enumerate}
        \item[(C1)] $\sigma(p) < \sigma(q)$
    \end{enumerate} 
\end{proposition}

\begin{lemma}\label{lem: criteria for equal permutation matrices}
    Given $a , b \in \nn$, and $\sigma , \omega \in S_\infty$, we have that $\sigma_{a , b} = \omega_{a , b}$ if the following conditions are satisfied.  
    \begin{enumerate}
        \item $r_{<a , b}(\sigma) = r_{<a , b}(\omega)$
        \item $r_{a , <b}(\sigma) = r_{a , <b}(\omega)$
        \item $r_{<a , <b}(\sigma) = r_{<a , <b}(\omega)$
        \item $r_{a , b}(\sigma) = r_{a , b}(\omega)$
    \end{enumerate} 
\end{lemma}

\begin{proposition}\label{prop: cover relations in the bruhat order}
    Let $\sigma \in S_\infty$ and $p, q \in \nn$ be such that $p < q$. Then $\sigma \lessdot_\text{Bruhat} \sigma \circ (p , q)$ if and only if the following conditions are satisfied. 
    \begin{enumerate}
        \item[(C1)] $\sigma(p) < \sigma(q)$
        \item[(C2)] For all $p < \ell < q$ we have that either $\sigma(\ell) < \sigma(p)$ or $\sigma(\ell) > \sigma(q)$.
    \end{enumerate}
\end{proposition}

\begin{proof}
First of all, suppose $\sigma \lessdot_\text{Bruhat} \sigma \circ (p , q)$. By Proposition \ref{prop: transposition going up in bruhat} we have $\sigma(p) < \sigma(q)$, so (C1) is satisfied. 

To show (C2), suppose that $p < \ell < q$ and that $\sigma(p) < \sigma(\ell) < \sigma(q)$. We claim that $\sigma <_\text{Bruhat} \sigma \circ (p , \ell) <_\text{Bruhat} \sigma \circ (p , q)$. The first inequality holds by Proposition \ref{prop: transposition going up in bruhat}. On the other hand, we have that $\sigma \circ (p , q) = \sigma \circ (p , \ell) \circ (\ell , q) \circ (p , \ell)$. Note that $[\sigma \circ (p , \ell)](\ell) = \sigma(p) < \sigma(q) = [\sigma \circ (p , \ell)](q)$, therefore by Proposition \ref{prop: transposition going up in bruhat} we have $\sigma \circ (p , \ell) <_\text{Bruhat} \sigma \circ (p , \ell) \circ (\ell , q)$. Similarly we have that $[\sigma \circ (p , \ell) \circ (\ell , q)](p) = \sigma(\ell) < \sigma(q) = [\sigma \circ (p , \ell) \circ (\ell , q)](\ell)$, so again by Proposition \ref{prop: transposition going up in bruhat} we have $\sigma \circ (p , \ell) \circ (\ell , q) <_\text{Bruhat} \sigma \circ (p , \ell) \circ (\ell , q) \circ (p , \ell) = \sigma \circ (p , q)$. So indeed we have $\sigma <_\text{Bruhat} \sigma \circ (p , \ell) <_\text{Bruhat} \sigma \circ (p , q)$ contradicting that $\sigma \lessdot_\text{Bruhat} \sigma \circ (p , q)$. Therefore (C2) is satisfied as well.

Now suppose $(p , q)$ satisfies (C1) and (C2). By Proposition \ref{prop: transposition going up in bruhat} we have $\sigma <_\text{Bruhat} \sigma \circ (p , q)$. Suppose that for some $\tau \in S_\infty$ we have $\sigma \leq_\text{Bruhat} \tau \leq_\text{Bruhat} \sigma \circ (p , q)$. By Proposition \ref{prop: region criterion for bruhat order}, for all $a , b \in \nn$ we have $r_{a , b}(\sigma) \geq r_{a , b}(\tau) \geq r_{a , b}(\sigma \circ (p , q))$. However by Lemma \ref{lem: ones in boxes after applying a transposition}, if $(\ast)$ $b < p$, $b \geq q$, $a < \sigma(p)$, or $a \geq \sigma(q)$ then $r_{a , b}(\sigma) = r_{ a , b}(\sigma \circ (p , q))$, implying that $r_{a , b}(\sigma) = r_{a , b}(\tau) = r_{ a , b}(\sigma \circ (p , q))$, and if $(\ast \ast)$ $\sigma(p) \leq a < \sigma(q)$ and $p \leq b < q$, then $r_{a , b}(\sigma) = r_{a , b}(\sigma \circ (p , q)) + 1$, so either $r_{a , b}(\tau) = r_{a , b}(\sigma)$ or $r_{a , b}(\tau) = r_{ a , b}(\sigma \circ (p , q))$.

First, we claim that $(\ast \ast \ast)$ $\sigma_{i , j} = \tau_{i , j}$ in the following (not disjoint) regions: (A) $i \in \nn$ and $j < p$, (B) $i < \sigma
(p)$ and $j \in \nn$, (C) $\sigma(q) < i$ and $j \in \nn$, (D) $i \in \nn$ and $q < j$. In any of the cases, since $\nn \times \nn$ is well-founded in the product order, there is a minimal $(i , j)$ in the relevant region such that $\sigma_{i , j} \neq \tau_{i , j}$. But $(\ast)$ implies that with $a = i$ and $b  = j$ the hypotheses of Lemma \ref{lem: criteria for equal permutation matrices} are satisfied, hence in fact $\sigma_{i , j} = \tau_{i , j}$, contradiction. 

Now we claim that $\sigma(p) \leq \tau(p) \leq \sigma(q)$. Indeed, we compute:  
\begin{equation*}
    r_{<\sigma(p) , < p}(\tau) \overset{(1)}{=} r_{<\sigma(p) , < p}(\sigma) \overset{(2)}{=} r_{<\sigma(p) , p}(\sigma) \overset{(3)}{=} r_{<\sigma(p) , p}(\tau)
\end{equation*}
Here (1) and (3) follow from $(\ast)$ and (2) follows because for all $1 \leq i < \sigma(p)$ we have $\sigma_{i , p} = 0$.  This implies that $\tau_{i , p} = 0$ for all $1 \leq i < \sigma(p)$, i.e. $\sigma(p) \leq \tau(p)$. Similarly we compute: 
\begin{equation*}
    r_{\sigma(q) , <p}(\tau) + 1 \overset{(1)}{=} r_{\sigma(q) , <p}(\sigma) + 1 \overset{(2)}{=} r_{\sigma(q) , p}(\sigma) \overset{(3)}{=} r_{\sigma(q) , p}(\tau)
\end{equation*}
Here again (1) and (3) follow from $(\ast)$ and (2) follows because $\sigma(p) < \sigma(q)$ by (C1) and $\sigma_{i , p} = 0$ unless $i = \sigma(p)$ in which case $\sigma_{i , p} = 1$. Therefore $\tau_{i , p} = 1$ for some $1 \leq i \leq \sigma(q)$, i.e. $\tau(p) \leq \sigma(q)$. Putting these two inequalities together yields the desired claim.  

Next we claim that in fact $\tau(p) \in \{ \sigma(p), \sigma(q) \}$. Suppose, for contradiction, that $\sigma(p) < \tau(p) < \sigma(q)$ and let $j = \sigma^{-1}(\tau(p))$, i.e. $\sigma_{\tau(p) , j} = 1$. Then we cannot have $j < p$, since by $(\ast \ast \ast)$ region (A), we have $\sigma_{i , j} = \tau_{i , j}$ for all $i \in \nn$ and $j < p$ . We also cannot have $j = p$, since $\sigma(p) < \tau(p)$ by hypothesis. Furthermore (C2), together with the fact that $\sigma(p) < \tau(p) < \sigma(q)$, implies that we cannot have $p < j < q$ either. Finally we cannot have $j = q$ since $\sigma(q) > \tau(p)$. Thus it must be that $j > q$. But then by $(\ast \ast \ast)$ region (D) we have that $\tau_{\tau(p) , j} = \sigma_{\tau(p) , j} = 1$, i.e. $\tau(j) = \tau(p)$, contradicting that $\tau$ is a bijection, so the claim holds.

Finally, we claim that $r_{[\sigma(p) , \sigma(q)],[p , q)}(\tau) = 1$. We first break up the region $[1 , \sigma(q)] \times [1 , q)$ into three disjoint reigons $[\sigma(p) , \sigma(q)] \times [1 , p)$, $[1 , \sigma(p)) \times [1 , q)$, and  $[\sigma(p) , \sigma(q)] \times [p , q)$, yielding the following. 
\begin{align*}
    r_{[\sigma(p) , \sigma(q)] , [1 , p)}(\tau) + r_{[1 , \sigma(p)) , [1 , q)}(\tau) + r_{[\sigma(p) , \sigma(q)],[p , q)}(\tau) 
    &= r_{[1 , \sigma(q)],[1 , q)}(\tau) 
    \\&\overset{(1)}{=} r_{[1 , \sigma(q)],[1 , q)}(\sigma) 
    \\&= r_{[\sigma(p) , \sigma(q)] , [1 , p)}(\sigma) + r_{[1 , \sigma(p)) , [1 , q)}(\sigma) + r_{[\sigma(p) , \sigma(q)],[p , q)}(\sigma)
\end{align*}
Here (1) follows from $(\ast)$. Since $\sigma(j) = \tau(j)$ for all $j < p$ we have $r_{[\sigma(p) , \sigma(q)] , [1 , p)}(\sigma) = r_{[\sigma(p) , \sigma(q)] , [1 , p)}(\tau)$. Also we have $r_{[1 , \sigma(p)) , [1 , q)}(\sigma) = r_{[1 , \sigma(p)) , [1 , q)}(\tau)$ again by $(\ast)$. Therefore subtraction yields $r_{[\sigma(p) , \sigma(q)],[p , q)}(\tau) = r_{[\sigma(p) , \sigma(q)],[p , q)}(\sigma)$. By (C2), we know that $r_{[\sigma(p) , \sigma(q)],[p , q)}(\sigma) = 1$, so the claim follows. 

Since $\tau(p) \in \{ \sigma(p) , \sigma(q) \}$, the previous claim implies that $\tau_{i , j} = 0$ for all $p \leq j < q$ and $\sigma(p) \leq i \leq \sigma(q)$ unless $j = p$ and either $i = \sigma(p)$ or $i = \sigma(q)$. 

By $(\ast \ast \ast)$ also we have that $\tau_{i , j} = \sigma_{i , j}$ in regions (A), (B), (C), and (D), so the only $(i , j)$ for which $\tau_{i , j}$ and $\sigma_{i , j}$ can differ is $(i , j) \in \{ (\sigma(p) , p), (\sigma(q), p) , (a , q) \mid \sigma(p) \leq a \leq \sigma(q) \}$, and exactly one of $\tau_{\sigma(p) , p}$ and $\tau_{\sigma(q) , p}$ must be $1$. Therefore $\tau(j) = \sigma(j)$ for all $j \neq p , q$, and so because $\tau$ is a bijection it must be that in the former case, $\tau(q) = \sigma(q)$ and in the latter case $\tau(q) = \sigma(p)$. Thus in the former case $\tau = \sigma$ and in the latter $\tau = \sigma \circ (p , q)$, as desired.
\end{proof}

\end{document}
